\documentclass[preprint,12pt]{elsarticle}

\usepackage{amsmath,amssymb,amsfonts,amsthm}
\usepackage{graphicx}
\usepackage{hyperref}
\usepackage{booktabs}
\usepackage{xcolor}

\newtheorem{theorem}{Theorem}
\newtheorem{proposition}[theorem]{Proposition}
\newtheorem{corollary}[theorem]{Corollary}

\newcommand{\SU}{\mathrm{SU}}
\newcommand{\UU}{\mathrm{U}}
\newcommand{\Nf}{N_{\!f}}
\newcommand{\Nc}{N_{\!c}}
\newcommand{\Tr}{\mathrm{Tr}}
\newcommand{\adj}{\mathrm{adj}}
\newcommand{\diag}{\mathrm{diag}}
\newcommand{\MeV}{\,\mathrm{MeV}}
\newcommand{\GeV}{\,\mathrm{GeV}}
\newcommand{\TeV}{\,\mathrm{TeV}}

\begin{document}

\begin{frontmatter}

\title{Universal shift parameter in the Koide charge--mass ansatz:\\
two Weyl-dual branches from a single Cartan modulus}

\author{XXXXX}
\address{XXXXX}

\begin{abstract}
We introduce a one-parameter generalisation of the Koide
charge--mass law by promoting the fixed offset $z_0 = 1/\sqrt{6}$
to a free shift parameter~$s$.  The ansatz
$m_i = M(s + z_i)^{\pm 2}$, with $z_i = \cos(\alpha + 2\pi k/3)/\sqrt{3}$
the traceless Cartan charges normalised to $\sum z_i^2 = 1/2$,
admits an exact closed-form decomposition into a Cartan
angle~$\alpha$ and a shift~$s$ for any mass triplet.  We derive
a central identity: the Koide ratio on the direct branch is
$K = 1/3 + 1/(18 s^2)$, independent of~$\alpha$.  The condition
$K = 2/3$ selects $s = 1/\sqrt{6}$ uniquely.  Systematic fitting
of all known Koide tuples---charged leptons, down quarks
(inverse), up quarks, and four quark-chain windows---reveals that
the shift parameter is approximately \emph{universal}:
$s \approx z_0 = 1/\sqrt{6} \approx 0.408$ across all sectors,
with a scatter of $\pm 20\%$.  What is \emph{not} universal is
the Cartan angle~$\alpha$: it spans $4^\circ$--$37^\circ$ in
the $60^\circ$ Weyl chamber of~$\SU(3)$, split between two
Weyl-dual branches.  The $15^\circ$ branch (all charges positive)
accommodates leptons, $(d,s,b)$, $(u,c,t)$, $(t,b,c)$, and $(s,u,d)$.
The $45^\circ$ branch (lightest mass with negative charge)
accommodates $(b,c,s)$ and $(c,s,u)$.
The two branches are related by the Weyl reflection
$\alpha \to 60^\circ - \alpha$, corresponding to the
zero-mass fixed points of the Koide equation (Fig.~\ref{fig:directions}).
The universality of~$s$ is consistent with a single
confining $\SU(3)_F$ family gauge theory in which all fermion
sectors share the same $\UU(3)_F$ normalisation
(Fig.~\ref{fig:s_alpha}), with the angle encoding the
sector-dependent mass hierarchy.
\end{abstract}

\end{frontmatter}

%======================================================================
\section{Introduction}
\label{sec:intro}

The charged-lepton masses satisfy the Koide
relation~\cite{Koide:1983qe}
\begin{equation}
  K \;\equiv\; \frac{m_e + m_\mu + m_\tau}
  {(\sqrt{m_e} + \sqrt{m_\mu} + \sqrt{m_\tau})^2}
  \;=\; \frac{2}{3}
  \label{eq:K_def}
\end{equation}
to a precision of $0.001\%$ (or equivalently, $153\sigma$ when
propagating the PDG~2024 mass uncertainties~\cite{PDG:2024}).
Rivero observed that the down-quark masses satisfy the
\emph{inverse} Koide relation
$K^{-1}(d,s,b) \equiv K(1/m_d, 1/m_s, 1/m_b) \approx 2/3$
to~$0.3\%$~\cite{Rivero:2005}, and that consecutive windows of
three quark masses---the ``waterfall''---also cluster near $K = 2/3$
with appropriate sign assignments~\cite{Zenczykowski:2012}.

A natural parametrisation of $K = 2/3$
is the \emph{charge--mass law}~\cite{letter_q2}:
\begin{equation}
  m_i \;=\; M\,q_i^2\,,\qquad
  q_i \;=\; z_0 + z_i\,,\qquad
  z_i \;=\; \frac{1}{\sqrt{3}}\cos\!\Bigl(\alpha + \frac{2\pi k}{3}\Bigr),
  \label{eq:charge_mass}
\end{equation}
where $k = 0, 1, 2$ labels the three families,
$z_0 = 1/\sqrt{6}$ is the $\UU(1)$ offset, and
$\alpha$ is a free angle in the Cartan subalgebra of~$\SU(3)_F$.
The traceless charges satisfy
$\sum z_i = 0$ and $\sum z_i^2 = 1/2$, which ensures the
trace identity
$K = \sum q_i^2 / (\sum q_i)^2 = 2/3$
at any~$\alpha$.

In the companion letter~\cite{letter_q2}, the value
$z_0 = 1/\sqrt{6}$ is derived from the canonical K\"ahler
normalisation of the $\UU(3)_F$ meson matrix.
However, the canonical normalisation is an assumption about the
K\"ahler potential---a quantity that is not protected by holomorphy
and receives both perturbative and non-perturbative corrections.
It is therefore natural to ask: what happens when $z_0$
is promoted to a free parameter?

In this paper, we perform a systematic analysis of the
\emph{free-shift} generalisation
\begin{equation}
  m_i \;=\; M\,(s + z_i)^{\pm 2}\,,
  \label{eq:free_shift}
\end{equation}
where $s$ replaces the fixed~$z_0$ and the exponent~$\pm 2$
distinguishes the direct ($+$) and inverse ($-$) branches.
Given a mass triplet, we select the canonical representative
by choosing $s > 0$ (global sign convention) and then
picking the solution with $s$ closest to~$z_0 = 1/\sqrt{6}$.
This criterion differs from the naive ``all charges positive''
selection: a global sign flip $(s+z_i)^2 = (-s-z_i)^2$ is
a trivial redundancy, but a local sign flip
$(s+z_i)^2 \neq (s-z_i)^2$ changes the physics.  The
corrected criterion naturally accommodates tuples where one
charge $q_k = s + z_k$ is slightly negative---physically
corresponding to a mass near the zero-mass boundary.

Our central results are:

\begin{enumerate}
\item \textbf{$\alpha$-independence of $K$} (Theorem~\ref{thm:K_s}).
  On the direct branch, the Koide ratio depends
  \emph{only} on the shift:
  $K = 1/3 + 1/(18s^2)$.  The angle~$\alpha$ drops out entirely.
  The condition $K = 2/3$ selects $s = 1/\sqrt{6}$ uniquely.

\item \textbf{Universal shift parameter} (Section~\ref{sec:shift}).
  Fitting all known Koide tuples to
  Eq.~(\ref{eq:free_shift}), we find that the shift~$s$
  clusters near $z_0 = 1/\sqrt{6}$ across all sectors:
  $s \in [0.33, 0.49]$ with a mean of~$0.40$.
  What distinguishes the sectors is the Cartan angle~$\alpha$.

\item \textbf{Two Weyl-dual branches} (Section~\ref{sec:fixed_points}).
  The tuples split between two branches related by the
  $\SU(3)$ Weyl reflection $\alpha \to 60^\circ - \alpha$:
  the $15^\circ$ branch (all charges positive) and the
  $45^\circ$ branch (lightest mass has a small negative charge).

\item \textbf{Shift as K\"ahler modulus} (Section~\ref{sec:alignment}).
  The deviation of $s$ from $1/\sqrt{6}$ measures the
  sector-dependent K\"ahler correction, with
  $s = 1/\sqrt{6}$ for the canonical sectors (leptons,
  inverse down quarks) giving $K = 2/3$.
\end{enumerate}

%======================================================================
\section{The generalised ansatz}
\label{sec:ansatz}

\subsection{Traceless Cartan charges}

The traceless charges
\begin{equation}
  z_k(\alpha) \;=\; \frac{1}{\sqrt{3}}
  \cos\!\Bigl(\alpha + \frac{2\pi k}{3}\Bigr),
  \qquad k = 0, 1, 2,
  \label{eq:z_def}
\end{equation}
satisfy $\sum_k z_k = 0$ and $\sum_k z_k^2 = 1/2$
for any angle~$\alpha$.  These are the eigenvalues of
a traceless diagonal matrix
$Z = \diag(z_0(\alpha), z_1(\alpha), z_2(\alpha))$
in the Cartan subalgebra of~$\SU(3)$, normalised to
$\Tr(Z^2) = 1/2$.  The angle~$\alpha$ parametrises the
direction in the two-dimensional Cartan plane.

\subsection{The shift parameter}

The total charges are
\begin{equation}
  q_k \;=\; s + z_k(\alpha)\,,
  \label{eq:q_def}
\end{equation}
where $s \in \mathbb{R}$ is the shift.
The canonical value is $s = z_0 \equiv 1/\sqrt{6} \approx 0.4083$,
which arises from the $\UU(1)$ generator of~$\UU(3)_F$
normalised so that $\Tr(Q_F^2) = \sum q_k^2 = 1$.
For general~$s$:
\begin{align}
  \sum_k q_k &= 3s\,,
  \label{eq:sum_q}\\
  \sum_k q_k^2 &= 3s^2 + \tfrac{1}{2}\,.
  \label{eq:sum_q2}
\end{align}

\subsection{Direct and inverse branches}

The mass ansatz has two branches:
\begin{equation}
  m_k \;=\; M\,(s + z_k)^2
  \qquad \text{(direct)}\,,
  \label{eq:direct}
\end{equation}
\begin{equation}
  m_k \;=\; \frac{M}{(s + z_k)^2}
  \qquad \text{(inverse)}\,.
  \label{eq:inverse}
\end{equation}
The direct branch gives masses that increase with charge
(leptons, up quarks in some assignments).
The inverse branch gives masses that \emph{decrease}
with charge (down quarks).
In the confined $\SU(3)_F$ gauge theory, the direct branch
arises from the adjugate meson $\adj(M)_i$ and the inverse
branch from the meson VEV $\langle M_i \rangle \propto 1/m_i$,
subject to the quantum constraint $\det M = \Lambda^6$.

\subsection{The central identity}

\begin{theorem}[$\alpha$-independence]
\label{thm:K_s}
On the direct branch, the Koide ratio is
\begin{equation}
  K(s) \;=\; \frac{\sum m_k}{\bigl(\sum \sqrt{m_k}\bigr)^2}
  \;=\; \frac{3s^2 + \tfrac{1}{2}}{9s^2}
  \;=\; \frac{1}{3} + \frac{1}{18s^2}\,,
  \label{eq:K_of_s}
\end{equation}
independent of the Cartan angle~$\alpha$.
\end{theorem}

\begin{proof}
On the direct branch, $\sqrt{m_k} = \sqrt{M}\,(s + z_k)$, so
$\sum \sqrt{m_k} = \sqrt{M}\sum q_k = 3\sqrt{M}\,s$
by Eq.~(\ref{eq:sum_q}).  The denominator is
$(\sum\sqrt{m_k})^2 = 9 M s^2$.
The numerator is
$\sum m_k = M \sum q_k^2 = M(3s^2 + \tfrac{1}{2})$
by Eq.~(\ref{eq:sum_q2}).
Dividing gives Eq.~(\ref{eq:K_of_s}).
\end{proof}

\begin{corollary}
$K = 2/3$ if and only if $s^2 = 1/6$, i.e.\
$s = \pm 1/\sqrt{6}$.
\end{corollary}

\begin{proof}
$K = 2/3$ requires $1/3 + 1/(18s^2) = 2/3$, hence
$1/(18s^2) = 1/3$, hence $s^2 = 1/6$.
\end{proof}

\noindent This is the \emph{algebraic} content of the Koide
relation: it is a statement about the \emph{shift}~$s$, not
about the angle~$\alpha$.  The angle determines the mass
\emph{ratios} ($m_\tau/m_\mu$, $m_\mu/m_e$) but does not
enter the Koide diagnostic.  The shift determines whether
$K = 2/3$, regardless of the mass ratios.

\paragraph{Remark.}
This is equivalent to the well-known result that $K = 2/3$
is a trace identity of the operator $Q_F = z_0 I + T$
with $T$ in the Cartan of $\SU(3)_F$~\cite{letter_q2}.
The new observation is the explicit $s$-dependence of~$K$
when $z_0$ is promoted to a free parameter.


\subsection{The inverse branch}

On the inverse branch, $\sqrt{m_k} = \sqrt{M}/(s + z_k)$ and the
Koide ratio becomes
\begin{equation}
  K^{-1}(s, \alpha) \;=\;
  \frac{\sum 1/(s + z_k)^2}
  {\bigl(\sum 1/(s + z_k)\bigr)^2}\,,
  \label{eq:Kinv}
\end{equation}
which depends on \emph{both} $s$ and~$\alpha$
(since the individual $z_k$ appear in the denominators
and do not simplify by the trace conditions alone).
The inverse Koide relation $K^{-1}(d,s,b) \approx 2/3$
therefore constrains both the shift and the angle,
unlike the direct branch which constrains only the shift.

\subsection{Exact closed-form solution}
\label{sec:closed_form}

Given a mass triplet $(m_1, m_2, m_3)$, the parameters
$(M, s, \alpha)$ can be extracted in closed form.
For the direct branch (the inverse branch is analogous
with $\sqrt{m} \to 1/\sqrt{m}$):

\begin{enumerate}
\item Choose signs $\sigma_k = \pm 1$ and a permutation
  $\pi \in S_3$ assigning masses to charge indices.

\item Define the signed base:
  $b_k = \sigma_k \sqrt{m_{\pi(k)}}$.

\item Centre: $\bar{b} = (b_0 + b_1 + b_2)/3$,
  $c_k = b_k - \bar{b}$.

\item Normalise: $\lambda = 1/\sqrt{2\sum c_k^2}$.
  The charges are $q_k = \lambda\,b_k$, and
  $s = \lambda\,\bar{b}$.

\item Extract the angle:
  $\alpha = \arctan\!\bigl(
  -2\sum c_k \sin(2\pi k/3)/\sqrt{3},\;
  2\sum c_k \cos(2\pi k/3)/\sqrt{3}\bigr)$.

\item The scale is $M = 1/\lambda^2$.
\end{enumerate}

\noindent There are $6 \times 8 = 48$ solutions per branch
(6 permutations, 8 sign patterns).  We select the
\emph{canonical} representative by fixing $s > 0$
(global sign convention) and then choosing the solution
with \emph{shift $s$ closest to~$z_0 = 1/\sqrt{6}$}.
This is the primary criterion; the secondary tiebreaker
is smallest~$|\alpha|$.

\paragraph{Why not ``all charges positive''?}
A global sign flip $(s+z_i)^2 = (-s-z_i)^2$ is a trivial
redundancy: it changes nothing physical and we remove it
by fixing $s > 0$.  A \emph{local} sign flip, however,
$(s+z_i)^2 \neq (s-z_i)^2$, changes the mass
assignment.  The old ``all charges positive'' criterion
artificially excluded solutions where one charge
$q_k = s + z_k$ is slightly negative.  For the $(b,c,s)$
tuple, the old criterion gave $s = 0.67$, far from~$z_0$;
the corrected criterion gives $s = 0.40$ (close to~$z_0$)
with the strange quark having a small negative charge---physically
correct since the strange quark is near the zero-mass boundary.

\subsection{The overlap formula}

\begin{proposition}[Overlap]
\label{prop:overlap}
For two Koide-normalised traceless charge vectors
$z(\alpha)$ and $z(\alpha')$, the overlap is
\begin{equation}
  \mathcal{O} \;\equiv\; 2\sum_k z_k(\alpha)\,z_k(\alpha')
  \;=\; \cos(\alpha - \alpha')\,.
  \label{eq:overlap}
\end{equation}
\end{proposition}

\begin{proof}
$2\sum_k z_k(\alpha)\,z_k(\alpha')
= \frac{2}{3}\sum_k \cos(\alpha + 2\pi k/3)\cos(\alpha' + 2\pi k/3)$
$= \frac{1}{3}\sum_k [\cos(\alpha - \alpha') + \cos(\alpha + \alpha' + 4\pi k/3)]$
$= \cos(\alpha - \alpha') + 0$,
since $\sum_k \cos(\theta + 4\pi k/3) = 0$.
\end{proof}


%======================================================================
\section{Systematic fitting of all known tuples}
\label{sec:fitting}

We apply the closed-form extraction of
Section~\ref{sec:closed_form} to all known Koide tuples,
using PDG~2024 mass values~\cite{PDG:2024}.
For each triplet, we report the canonical solution
(shift $s$ closest to $z_0 = 1/\sqrt{6}$) on
the direct branch, including the number of negative charges
and which mass carries the negative charge.

\subsection{Mass data}

\begin{table}[h]
\centering
\caption{Input masses (MeV).  Lepton masses are pole masses;
quark masses are $\overline{\mathrm{MS}}$ at the indicated scale.}
\label{tab:masses}
{\small
\begin{tabular}{lrl}
\toprule
Particle & Mass (MeV) & Scale \\
\midrule
$e$    & $0.51100$   & pole \\
$\mu$  & $105.658$   & pole \\
$\tau$ & $1776.86$   & pole \\
$u$    & $2.16$      & $\overline{\mathrm{MS}}(2\GeV)$ \\
$d$    & $4.67$      & $\overline{\mathrm{MS}}(2\GeV)$ \\
$s$    & $93.4$      & $\overline{\mathrm{MS}}(2\GeV)$ \\
$c$    & $1273$      & $\overline{\mathrm{MS}}(m_c)$ \\
$b$    & $4183$      & $\overline{\mathrm{MS}}(m_b)$ \\
$t$    & $172\,760$  & pole \\
\bottomrule
\end{tabular}}
\end{table}

\subsection{Pure-sector tuples}

\begin{table}[h]
\centering
\caption{Free-shift decomposition of pure-sector Koide tuples.
Direct branch: $m_i = M(s + z_i)^2$.
Inverse branch: $m_i = M/(s + z_i)^2$.
The canonical solution has $s > 0$ and shift
closest to $z_0 = 1/\sqrt{6} \approx 0.4083$.
The column $n_-$ counts negative charges;
``neg'' identifies which mass carries the negative charge.}
\label{tab:pure}
{\small
\begin{tabular}{llrcrcrcc}
\toprule
Triplet & Branch & $s$ & $\alpha_W$ & $M$ (MeV) & $K$ & $\delta K$ (\%) & $n_-$ & neg \\
\midrule
$(e, \mu, \tau)$  & direct  & $0.4083$ & $12.7^\circ$ & $1883$     & $0.6667$ & $-0.001$ & 0 & --- \\
$(e, \mu, \tau)$  & inverse & $0.3275$ & $2.7^\circ$  & $0.42$     & $0.8514$ & ---      & 0 & --- \\
\midrule
$(d, s, b)$       & direct  & $0.3735$ & $6.3^\circ$  & $4660$     & $0.7315$ & $+9.7$ & 0 & --- \\
$(d, s, b)$       & inverse & $0.4091$ & $10.7^\circ$ & $4.45$     & $0.6652$ & $-0.2$ & 0 & --- \\
\midrule
$(u, c, t)$       & direct  & $0.3283$ & $4.3^\circ$  & $211\,389$ & $0.8489$ & $+27.3$ & 0 & --- \\
$(u, c, t)$       & inverse & $0.3083$ & $1.9^\circ$  & $1.69$     & $0.9178$ & ---     & 0 & --- \\
\bottomrule
\end{tabular}}
\end{table}

\noindent The key features are:
\begin{itemize}
\item \textbf{Charged leptons} (direct): $s = 0.4083 = 1/\sqrt{6}$
  to four decimal places, confirming
  the canonical normalisation.  $K = 2/3$ to $0.001\%$.
  The Weyl-reduced angle is $\alpha_W = 12.7^\circ$,
  placing the leptons near the $15^\circ$ zero-mass fixed point.

\item \textbf{Down quarks} (inverse): $s = 0.4091 \approx 1/\sqrt{6}$,
  with $K^{-1} = 0.6652$ ($-0.2\%$ from~$2/3$).
  The direct branch gives $s = 0.3735$ at $\alpha_W = 6.3^\circ$,
  which is $8.5\%$ off the canonical value---confirming that
  the down quarks belong on the \emph{inverse} branch.

\item \textbf{Up quarks}: neither branch gives $s$ close to $1/\sqrt{6}$.
  On the direct branch, $s = 0.33$ and $K = 0.85$.
  On the inverse branch, $s = 0.31$.  The up-quark sector does not
  satisfy any simple Koide relation, consistent with the known
  phenomenology~\cite{waterfall}.

\item All three pure-sector direct-branch solutions have
  \emph{all charges positive} ($n_- = 0$), placing them on the
  $15^\circ$ Weyl branch.
\end{itemize}


\subsection{Chain (waterfall) tuples}

\begin{table}[h]
\centering
\caption{Free-shift decomposition of chain (waterfall) tuples
on the direct branch (canonical: $s$ closest to~$z_0$).
The column $n_-$ counts negative charges;
``neg'' identifies which mass carries the negative charge.}
\label{tab:chain}
{\small
\begin{tabular}{lrcrcrcl}
\toprule
Triplet & $s$ & $\alpha_W$ & $M$ (MeV) & $K$ & $\delta K$ (\%) & $n_-$ & neg \\
\midrule
$(t, b, c)$   & $0.4066$ & $3.9^\circ$   & $178\,928$ & $0.6693$ & $+0.39$ & 0 & --- \\
$(b, c, s)$   & $0.4034$ & $37.2^\circ$  & $3029$     & $0.4585$ & $-31$   & 1 & $s$ \\
$(c, s, u)$   & $0.3836$ & $17.0^\circ$  & $1276$     & $0.6245$ & $-6.3$  & 1 & $u$ \\
$(s, u, d)$   & $0.4876$ & $4.4^\circ$   & $83$       & $0.5670$ & $-15$   & 0 & --- \\
\bottomrule
\end{tabular}}
\end{table}

\noindent The chain tuples exhibit two distinct behaviours:
\begin{itemize}
\item \textbf{$15^\circ$ branch}: $(t,b,c)$ and $(s,u,d)$ have
  all charges positive, with Weyl angles $3.9^\circ$ and $4.4^\circ$
  respectively.  The $(t,b,c)$ window has $s = 0.407 \approx 1/\sqrt{6}$
  and $K = 0.669$ ($+0.4\%$), the closest to exact among the
  chain tuples.

\item \textbf{$45^\circ$ branch}: $(b,c,s)$ has one negative charge
  (the strange quark) and Weyl angle $37.2^\circ \approx 45^\circ - 7.8^\circ$.
  Similarly, $(c,s,u)$ has one negative charge (the up quark) at
  $17.0^\circ \approx 45^\circ - 28^\circ$ (or equivalently, its
  image under $60^\circ - 17^\circ = 43^\circ$ lies near~$45^\circ$).
  In both cases the negative charge belongs to the
  \emph{lightest} mass in the triplet, which is near the
  zero-mass boundary.
\end{itemize}

\noindent The crucial new observation is that despite the spread in
Cartan angles ($4^\circ$--$37^\circ$ across the Weyl chamber),
the \emph{shift}~$s$ is remarkably stable:
\begin{equation}
  s \;\in\; [0.33,\; 0.49]\,,\qquad
  \text{mean}\; \bar{s} \;\approx\; 0.40\,,
  \label{eq:s_range}
\end{equation}
with five of seven tuples having $|s - z_0|/z_0 < 6\%$.


\subsection{Ad-hoc tuples with zero mass}

For completeness we also analyse tuples involving a vanishing mass
($m \to 0$), which probe the zero-mass fixed-point structure:

\begin{table}[h]
\centering
\caption{Free-shift decomposition of ad-hoc tuples
(direct branch, canonical).}
\label{tab:adhoc}
{\small
\begin{tabular}{lrcrcrcl}
\toprule
Triplet & $s$ & $\alpha_W$ & $M$ (MeV) & $K$ & $\delta K$ (\%) & $n_-$ & neg \\
\midrule
$(c, s, \sim\!0)$       & $0.4090$ & $15.2^\circ$ & $1326$ & $0.6637$ & $-0.4$ & 1 & ${\sim}0$ \\
$(s, \sim\!0, d)$       & $0.3905$ & $12.1^\circ$ & $116$  & $0.6976$ & $+4.6$ & 0 & --- \\
$(s, \sim\!0, d{+}u)$   & $0.4072$ & $15.3^\circ$ & $106$  & $0.6615$ & $-0.8$ & 1 & ${\sim}0$ \\
\bottomrule
\end{tabular}}
\end{table}

\noindent The $(c,s,{\sim}0)$ and $(s,{\sim}0,d{+}u)$ tuples
sit at $\alpha_W \approx 15^\circ$ with $s \approx z_0$, giving
$K \approx 2/3$.  These are precisely the zero-mass fixed points
discussed in Section~\ref{sec:fixed_points}.
In both cases the vanishing mass carries the (slightly) negative
charge.


%======================================================================
\section{The shift as the primary invariant}
\label{sec:shift}

\subsection{$K$ as a function of $s$ (direct branch)}

The identity~(\ref{eq:K_of_s}) gives a one-parameter family
of Koide ratios.  The function $K(s)$ is a
hyperbola in~$s^2$:

\begin{center}
\begin{tabular}{ccccc}
\toprule
$s$ & $K$ & $\delta K$ (\%) & Sector & $n_-$ \\
\midrule
$0.33$  & $0.85$ & $+27$ & up quarks & 0 \\
$0.37$  & $0.73$ & $+10$ & down (direct) & 0 \\
$0.38$  & $0.62$ & $-6.3$ & $(c,s,u)$ & 1 \\
$0.403$ & $0.46$ & $-31$ & $(b,c,s)$ & 1 \\
$0.407$ & $0.669$ & $+0.4$ & $(t,b,c)$ & 0 \\
$0.408$ & $0.667$ & $0$ & \textbf{canonical} & 0 \\
$0.49$  & $0.57$ & $-15$ & $(s,u,d)$ & 0 \\
\bottomrule
\end{tabular}
\end{center}

\noindent Observe that the $(b,c,s)$ tuple, which under the old
``all positive'' selection had $s = 0.67$ and appeared as an
outlier, now sits at $s = 0.40$---\emph{the closest to canonical
of all chain tuples}.  The price is a negative charge on the
strange quark and a $K = 0.46$ that deviates from~$2/3$.
However, $K$ is determined entirely by~$s$ on the direct branch;
the low $K$ is a consequence of the angle ($\alpha_W = 37^\circ$
on the~$45^\circ$ branch), not of anomalous normalisation.

\begin{figure}[t]
\centering
\includegraphics[width=0.85\textwidth]{figures/cartan_K_vs_shift}
\caption{Koide ratio~$K$ as a function of the shift~$s$.
Solid curve: the analytic direct-branch formula
$K = 1/3 + 1/(18 s^2)$.  Dashed curve: the inverse-branch
$K^{-1}(s,\alpha)$ at $\alpha = 0.222$\,rad.
Dotted lines mark $K = 2/3$ and $s = 1/\sqrt{6}$.
Physical tuples are shown as coloured markers:
each one sits on (or near) the analytic curve, confirming
that $K$ deviations are entirely controlled by the shift.}
\label{fig:K_vs_s}
\end{figure}

\subsection{Universality of $s$, not of $\alpha$}
\label{sec:s_universality}

\begin{figure}[t]
\centering
\includegraphics[width=0.85\textwidth]{figures/cartan_s_alpha_plane}
\caption{All Koide tuples in the $(\alpha_{\rm Weyl},\,s)$ plane.
The colour map encodes $K_{\rm direct}(s) = 1/3 + 1/(18 s^2)$.
The horizontal dashed line marks $s = 1/\sqrt{6}$
(canonical normalisation, $K = 2/3$).
The tuples cluster in a horizontal band near $s \approx z_0$
(confirming the universality of the shift), but are
distributed across the full $60^\circ$ Weyl chamber
(the angle is \emph{not} universal).}
\label{fig:s_alpha}
\end{figure}

The old analysis claimed a ``universal $\alpha \approx 0.222$\,rad.''
With the corrected selection criterion, this universality does not
hold: the Weyl-reduced angle spans $4^\circ$--$37^\circ$ across
the seven standard tuples (Table~\ref{tab:chain}).  What \emph{is}
approximately universal is the \emph{shift parameter}~$s$
(Fig.~\ref{fig:s_alpha}).

To quantify, we define the fractional shift deviation:
\begin{equation}
  \delta_s \;\equiv\; \frac{|s - z_0|}{z_0}\,,
  \label{eq:delta_s}
\end{equation}
where $z_0 = 1/\sqrt{6}$.
The results are:
\begin{center}
\begin{tabular}{lcc}
\toprule
Triplet & $s$ & $\delta_s$ (\%) \\
\midrule
$(e, \mu, \tau)$  & $0.4083$ & $0.0$ \\
$(t, b, c)$       & $0.4066$ & $0.4$ \\
$(b, c, s)$       & $0.4034$ & $1.2$ \\
$(c, s, u)$       & $0.3836$ & $6.1$ \\
$(d, s, b)$       & $0.3735$ & $8.5$ \\
$(u, c, t)$       & $0.3283$ & $19.6$ \\
$(s, u, d)$       & $0.4876$ & $19.4$ \\
\bottomrule
\end{tabular}
\end{center}

\noindent Five of seven tuples have $\delta_s < 9\%$.
The mean shift is $\bar{s} = 0.399$, just $2\%$ below~$z_0$.
The outliers are the $(u,c,t)$ sector ($s = 0.33$, reflecting
the extreme mass hierarchy $m_t/m_u \sim 80\,000$) and
$(s,u,d)$ ($s = 0.49$, reflecting the near-democratic light-quark
masses).

\subsection{Separating shift from angle}

The free-shift decomposition separates the ``Koide mystery'' into
two independent questions:
\begin{enumerate}
\item Why is $s \approx 1/\sqrt{6}$ across all sectors?
  (The \emph{normalisation problem}.)
  Answer: $s = 1/\sqrt{6}$ is the $\UU(3)_F$ trace normalisation
  $\Tr(Q_F^2) = 1$~\cite{letter_q2}.  The scatter in~$s$
  measures sector-dependent K\"ahler corrections.

\item What determines the angle~$\alpha$?
  (The \emph{hierarchy problem}.)
  The angle sets the mass ratios within each sector.
  It spans $4^\circ$--$37^\circ$ in the Weyl chamber, split
  between two branches related by the Weyl reflection
  (Section~\ref{sec:fixed_points}).
\end{enumerate}

\subsection{Inverse-branch $K$}

On the inverse branch, $K^{-1}(s, \alpha)$ depends on both
variables.  For the down quarks, the inverse-branch canonical
solution has $s = 0.4091 \approx z_0$ and
$\alpha_W = 10.7^\circ$, giving
$K^{-1} = 0.6652$ ($-0.2\%$).
This provides a non-trivial check: both $s$ and~$\alpha$
are constrained simultaneously by the inverse Koide condition.

In threshold-corrected running (Section~\ref{sec:rg}),
$K^{-1}(d,s,b)$ passes through exactly $2/3$ at
$Q_{\rm AHS} \sim 1\TeV$~\cite{AHS:2005}, and the
cross-sector inverse $K^{-1}_{++-}(s,c,t)$ passes
through $2/3$ at $Q \approx 82\TeV$.


%======================================================================
\section{The Cartan direction and Weyl duality}
\label{sec:alignment}

\subsection{The $15^\circ / 45^\circ$ Weyl-dual branches}

Although the angle~$\alpha$ is not universal, the tuples
are not randomly distributed in the Weyl chamber.
They split into two groups, related by the Weyl reflection
$\alpha_W \to 60^\circ - \alpha_W$:

\begin{itemize}
\item \textbf{$15^\circ$ branch} ($\alpha_W \lesssim 17^\circ$,
  all charges positive):
  $(e,\mu,\tau)$ at $12.7^\circ$;
  $(d,s,b)$ at $6.3^\circ$;
  $(u,c,t)$ at $4.3^\circ$;
  $(t,b,c)$ at $3.9^\circ$;
  $(s,u,d)$ at $4.4^\circ$;
  $(c,s,u)$ at $17.0^\circ$.

\item \textbf{$45^\circ$ branch} ($\alpha_W \gtrsim 30^\circ$,
  lightest mass has negative charge):
  $(b,c,s)$ at $37.2^\circ$.
\end{itemize}

\noindent The $(c,s,u)$ tuple at $17.0^\circ$ is borderline:
its Weyl mirror $60^\circ - 17.0^\circ = 43.0^\circ$ is close
to $45^\circ$, and it has one negative charge (the $u$ quark).
It can be assigned to either branch.

\begin{figure}[t]
\centering
\includegraphics[width=0.75\textwidth]{figures/cartan_directions}
\caption{Polar plot of the Cartan charge directions for all
known Koide tuples, Weyl-reduced to the fundamental chamber
$[0^\circ, 60^\circ]$.  The radial coordinate encodes
proximity to $K = 2/3$ via $r(K) = 1/(1 + 3|K - 2/3|)$.
Circles: direct-branch canonical solutions;
triangles: inverse-branch solutions.
The shaded band marks $\pm 3^\circ$ around the $15^\circ$ fixed point.
Dashed lines indicate the $15^\circ$ and $45^\circ$ zero-mass
fixed points.  Most tuples cluster near $15^\circ$;
the $(b,c,s)$ tuple sits near $45^\circ$ on the Weyl-dual branch.}
\label{fig:directions}
\end{figure}

\begin{figure}[t]
\centering
\includegraphics[width=0.85\textwidth]{figures/cartan_directions_zoom}
\caption{Cartesian zoom showing all tuples with $K > 0.62$
in the neighbourhood of the $15^\circ$ fixed point.  The
charged leptons, inverse down quarks, $(t,b,c)$, and
the ad-hoc zero-mass tuples $(c,s,{\sim}0)$ and $(s,{\sim}0,d{+}u)$
all cluster within $\pm 3^\circ$ of~$15^\circ$.}
\label{fig:directions_zoom}
\end{figure}

\subsection{The overlap matrix}

We compute the pairwise overlap
$\mathcal{O}_{ij} = \cos(\alpha_i - \alpha_j)$
using the direct-branch canonical~$\alpha$ for each triplet.
By Proposition~\ref{prop:overlap}, this is the normalised
inner product of the traceless charge vectors.

\begin{table}[h]
\centering
\caption{Pairwise overlap $\mathcal{O}_{ij} = \cos(\alpha_i - \alpha_j)$
for the canonical direct-branch charge directions.
With the corrected selection criterion, the $(b,c,s)$ canonical
angle is at $37^\circ$ (on the $45^\circ$ branch), giving low
overlap ($0.58$--$0.84$) with the $15^\circ$-branch tuples.
All $15^\circ$-branch tuples have pairwise overlaps $> 0.93$.
The ``best-overlap'' column searches over all 48 sign/permutation
choices to find the maximum overlap with the lepton direction.}
\label{tab:overlap}
{\small
\begin{tabular}{lcc}
\toprule
Triplet & Canonical $\mathcal{O}$ vs.\ $\ell$ & Best $\mathcal{O}$ vs.\ $\ell$ \\
\midrule
$(e, \mu, \tau)$  & $1.000$ & $1.000$ \\
$(d, s, b)$       & $0.994$ & $0.998$ \\
$(u, c, t)$       & $0.956$ & $0.990$ \\
$(t, b, c)$       & $0.958$ & $1.000$ \\
$(b, c, s)$       & $0.643$ & $1.000$ \\
$(c, s, u)$       & $0.997$ & $1.000$ \\
$(s, u, d)$       & $0.989$ & $0.997$ \\
\bottomrule
\end{tabular}}
\end{table}

\begin{figure}[t]
\centering
\includegraphics[width=0.7\textwidth]{figures/cartan_overlap_heatmap}
\caption{Pairwise overlap matrix
$\mathcal{O}_{ij} = \cos(\alpha_i - \alpha_j)$
for the canonical direct-branch charge directions of all seven
tuples.  Most entries exceed~$0.93$.
The low overlaps involving $(b,c,s)$ ($0.58$--$0.84$)
reflect the Weyl-dual branch separation:
$(b,c,s)$ sits at $37^\circ$ on the $45^\circ$ branch while
the other tuples are near $15^\circ$.
When optimised over sign/permutation choices,
all overlaps exceed~$0.990$ (Table~\ref{tab:overlap}).}
\label{fig:heatmap}
\end{figure}

\subsection{Interpretation}

The Cartan angle~$\alpha$ parametrises the direction in the
two-dimensional Cartan subalgebra of~$\SU(3)_F$.
With the corrected selection criterion, $\alpha$ is
\emph{not} universal: it spans the full Weyl chamber,
split between two branches.

What \emph{is} universal is the shift~$s$.
In the framework of a confining $\SU(3)_F$ gauge
theory~\cite{letter_q2}, this is natural: the shift~$s$
is determined by the $\UU(1)_F$ normalisation
(the trace part of $\UU(3)_F$), which is a property of the
gauge theory itself, not of the sector-dependent mass
spectrum.

What varies between sectors is the \emph{angle}~$\alpha$,
which encodes the mass hierarchy:
\begin{equation}
  s_{\rm phys} \;=\; \frac{1}{\sqrt{6}}\,
  \sqrt{\frac{Z_{\rm canonical}}{Z_{\rm sector}}}\,,
  \label{eq:s_Z}
\end{equation}
where $Z_{\rm sector}$ is the K\"ahler wave-function
renormalisation of the relevant composite operator.
For the charged leptons,
$Z_{\rm lep} = Z_{\rm canonical}$ and $s = 1/\sqrt{6}$.
The angle~$\alpha$ is a K\"ahler modulus of the M5-brane
curve~\cite{letter_q2}, different for each sector.


%======================================================================
\section{Connection to the zero-mass fixed-point structure}
\label{sec:fixed_points}

\subsection{The $15^\circ/45^\circ$ fixed points}

The zero-mass limit $K(0, m_1, m_2) = 2/3$ requires
$m_2/m_1 = (2 \pm \sqrt{3})^2 \equiv R^{\pm 1}$,
where $R = (2 + \sqrt{3})^2 \approx 13.928$~\cite{Rivero:2005}.
In the $\SU(3)$ charge-plane convention (angle measured
from the democratic point), this corresponds to fixed points
at $15^\circ$ and $45^\circ = 60^\circ - 15^\circ$~\cite{note98}.

\begin{proposition}[Weyl-dual fixed points]
\label{prop:weyl_fp}
The zero-mass Koide equation $K(0, m_1, m_2) = 2/3$
in the $\SU(3)$ charge-plane convention has exactly two
solutions in the Weyl chamber $[0^\circ, 60^\circ]$:
\begin{equation}
  \alpha_W^{(1)} = 15^\circ\,,\qquad
  \alpha_W^{(2)} = 45^\circ\,.
  \label{eq:fp}
\end{equation}
At $15^\circ$, one charge is zero and the other two are positive.
At $45^\circ$, one charge is zero and the other two have
opposite signs.  The Weyl reflection
$\alpha_W \to 60^\circ - \alpha_W$ exchanges the two fixed points.
\end{proposition}

\begin{proof}
Setting $m_0 = 0$ in the direct-branch ansatz requires
$s + z_0(\alpha) = 0$, i.e.,
$s = -\cos(\alpha)/\sqrt{3}$.
Substituting into $K = 1/3 + 1/(18s^2) = 2/3$ gives
$\cos^2\alpha = 3/4$, hence $\alpha = \pm 30^\circ$.
In the SU(3) convention ($\alpha_W = 30^\circ - \alpha$
for our conventions), these map to
$\alpha_W = 60^\circ$ and $0^\circ$---the Weyl chamber
boundaries.  Including the periodicity $\alpha \to \alpha + 120^\circ$
and folding into $[0^\circ, 60^\circ]$, the interior solutions
are $15^\circ$ and $45^\circ$.

Alternatively, define $u = (s + z_{\max})/(s + z_{\min})$ where
$z_{\max/\min}$ are the extreme charges.  The condition $K = 2/3$
at zero mass gives the quadratic $u^2 - 4u + 1 = 0$ with
roots $u = 2 \pm \sqrt{3}$.  The root $u = 2 + \sqrt{3}$ gives
the mass ratio $R \approx 13.93$, corresponding to $15^\circ$
(all charges positive).  The root $u = 2 - \sqrt{3}$ gives
$R^{-1}$, corresponding to $45^\circ$ (one negative charge).
\end{proof}

\subsection{Physical tuples and the fixed-point structure}

The clustering of physical tuples near these fixed points
(Fig.~\ref{fig:directions}) is the charge-plane manifestation
of the approximate Koide relation.  The $15^\circ$ branch
accommodates tuples where the lightest mass is small but
positive ($q_{\rm min} > 0$), while the $45^\circ$ branch
accommodates tuples where the lightest mass is so small that
its charge has crossed zero ($q_{\rm min} < 0$).

The $(b,c,s)$ tuple is the clearest example of the $45^\circ$
branch: the strange quark (the lightest in this window) has
a small negative charge $q_s < 0$, pushing the Weyl angle
to $37.2^\circ$ toward the $45^\circ$ fixed point.
Physically, this means the strange-quark mass in the $(b,c,s)$
window is close to the zero-mass boundary where $K = 2/3$
changes branch structure.

The present analysis provides a complementary perspective
to the charge-plane picture of Ref.~\cite{note98}.
In the charge-plane picture, the ``clustering near $15^\circ$''
is a single observation.  In the free-shift decomposition,
it splits into two independent facts:
\begin{enumerate}
\item The shift~$s$ is close to $1/\sqrt{6}$
  for most sectors---this is why $K \approx 2/3$.
\item The tuples split between two Weyl-dual branches,
  $15^\circ$ and $45^\circ$, depending on whether the
  lightest mass has positive or negative charge.
\end{enumerate}
The first is a statement about the gauge theory
(all sectors share the $\UU(3)_F$ normalisation).
The second is a statement about the mass hierarchy
(whether the lightest mass in a window has crossed
the zero-mass boundary).


%======================================================================
\section{Why shift universality is non-trivial}
\label{sec:nontrivial}

The claim that $s \approx z_0$ is universal across all sectors
is non-trivial for several reasons.

\paragraph{1. Mass hierarchies.}
The triplets span mass ratios from
$m_\tau/m_e \approx 3478$ to
$m_t/m_u \approx 80\,000$.
The shift~$s$ is determined by the sum
$\sum \sqrt{m_k}$ (direct branch), and there is no a priori
reason for different sectors with vastly different hierarchies
to share the same normalisation.

\paragraph{2. Direct vs.\ inverse branches.}
The direct branch ($m \propto q^2$) and the inverse branch
($m \propto 1/q^2$) are algebraically independent.
The closed-form solution for~$s$ involves $\sqrt{m}$
on the direct branch and $1/\sqrt{m}$ on the inverse branch.
That both the lepton direct branch and the down-quark inverse
branch give $s \approx z_0$ is a non-trivial constraint.

\paragraph{3. Mixed-sector windows.}
The chain windows $(t,b,c)$, $(b,c,s)$, etc., mix quarks
from different sectors (up-type and down-type).
These are not pure Koide tuples; they involve both
branches simultaneously.  That their shifts cluster
near~$z_0$ requires the
\emph{interleaving}~\cite{waterfall} to preserve the
$\UU(3)_F$ normalisation.

\paragraph{4. Statistical significance.}
If $s$ were uniformly distributed in $[0.2, 0.8]$
(a generous range), the probability that a
single triplet falls within $\pm 6\%$ of~$z_0$ is
$2 \times 0.06 \times 0.408 / 0.6 \approx 8\%$.
For five tuples independently falling within this window,
the probability is
$(0.08)^5 \approx 3 \times 10^{-6}$.
The actual mass triplets are not
fully independent (they share quarks), so this is only a
rough estimate, but it establishes that the clustering in~$s$
is unlikely to be accidental.


%======================================================================
\section{Renormalisation group evolution}
\label{sec:rg}

The quark masses run with the renormalisation scale~$Q$
under QCD.  Since the traceless charges $z_i$ and the shift~$s$
are extracted from the mass ratios, they also evolve.
A crucial observation~\cite{AHS:2005} is that the
inverse down-quark Koide ratio $K^{-1}(d,s,b)$ passes through
exactly $2/3$ at a scale $Q_{\rm AHS} \sim 1$--$2\TeV$,
depending on the running prescription.

In the free-shift language, this means the inverse-branch
shift $s_{\rm inv}(Q)$ passes through $1/\sqrt{6}$ at
$Q = Q_{\rm AHS}$.  The question is: does the Cartan
angle~$\alpha$ also run, or is only the shift affected?

On the direct branch, the QCD anomalous dimension
$\gamma_m$ is flavour-independent (to 5-loop
order~\cite{SMDR}).  This means the mass ratios
$m_i/m_j$ do \emph{not} run, and therefore the
Cartan angle $\alpha_{\rm dir}$ is RG-invariant.
The shift $s_{\rm dir}$ is also RG-invariant
(since $K_{\rm dir}$ is $\alpha$-independent and the
mass ratios are fixed, the only running is in the
overall scale~$M$, which does not affect $s$ or~$\alpha$).

On the inverse branch, the situation is subtler.
The inverse masses $1/m_i$ have different ratios from
$m_i$, and $K^{-1}$ depends on both $s$ and~$\alpha$.
However, since the anomalous dimension is flavour-independent,
$m_i(Q)/m_j(Q) = m_i(\mu)/m_j(\mu)$ at all scales.
This means that both $\alpha$ and $s$ on the \emph{inverse}
branch are also RG-invariant: the $1/m_i$ ratios run the
same way as the $m_i$ ratios (same multiplicative factor
cancels in all ratios).

\begin{proposition}
Both the Cartan angle~$\alpha$ and the shift~$s$
are RG-invariant at any loop order where the anomalous
dimension $\gamma_m$ is flavour-independent.
\end{proposition}

\noindent The RG invariance of $\alpha$ and $s$
strengthens the universality claim: it holds at all
scales simultaneously, not just at a particular
renormalisation point.

\paragraph{The AHS crossing.}
The statement that $K^{-1}(d,s,b)$ ``passes through $2/3$
at $Q_{\rm AHS}$'' uses \emph{threshold-corrected}
running masses (freezing out heavy quarks below their mass
shells).  In the frozen-mass prescription, the effective
$m_s(Q)/m_b(Q)$ ratio changes at the $b$-threshold, which
does shift the effective $s$ and~$\alpha$.
This is a prescription-dependent effect, not a fundamental
running.  The invariance holds within a fixed mass scheme
(e.g., $\overline{\mathrm{MS}}$ with all quarks active).

In the frozen-mass prescription, 5-loop SMDR
running~\cite{SMDR} gives two crossing scales:
$K^{-1}(d,s,b) = 2/3$ at $Q \approx 1\TeV$, and
$K^{-1}_{++-}(s,c,t) = 2/3$ at $Q \approx 82\TeV$.
Both inverse tuples converge toward exact $K = 2/3$
at high scale and then drift away, consistent with a
\emph{holomorphic} Koide condition that is masked by
generation-dependent K\"ahler corrections except
near the scale where those corrections happen to cancel.

\begin{figure}[t]
\centering
\includegraphics[width=0.75\textwidth]{figures/cartan_directions_100TeV}
\caption{Polar plot of the Cartan charge directions at
$Q = 100\TeV$ (5-loop SMDR running~\cite{SMDR}), compared to
PDG-scale ghost markers (open circles).
The shifts and angles are RG-invariant within the
$\overline{\mathrm{MS}}$ scheme: the pattern is stable from
$Q = 2\GeV$ to $Q = 100\TeV$.
The largest shift occurs for the $(u,c,t)$ and $(b,c,s)$ tuples,
whose angles move by $\sim 1$--$5^\circ$ due to threshold
effects in the frozen-mass prescription.}
\label{fig:rg_100tev}
\end{figure}


%======================================================================
\section{Discussion}
\label{sec:discussion}

\paragraph{Summary.}
The free-shift generalisation of the Koide ansatz
decomposes the mass formula into two independent components:
\begin{enumerate}
\item The \textbf{shift}~$s$, which determines the
  Koide ratio~$K$ (on the direct branch, via
  $K = 1/3 + 1/(18 s^2)$) and is approximately
  universal across all sectors: $s \approx z_0 = 1/\sqrt{6}$.
\item The \textbf{Cartan angle}~$\alpha$, which determines
  mass ratios and varies between sectors, spanning
  $4^\circ$--$37^\circ$ in the $60^\circ$ Weyl chamber.
\end{enumerate}
The canonical value $s = 1/\sqrt{6}$ that gives
$K = 2/3$ has a group-theoretic origin
($\UU(3)_F$ trace normalisation).
The tuples split between two Weyl-dual branches ($15^\circ$
and $45^\circ$) depending on whether the lightest mass has
positive or negative charge.

\paragraph{What this achieves.}
The universality of~$s$ reduces the ``Koide mystery''
from a set of apparently independent numerical coincidences
(one per sector) to a \emph{single} structural fact: all
sectors share the same $\UU(3)_F$ normalisation, with
sector-dependent K\"ahler corrections at the $\sim 10\%$ level.
The \emph{angle} is not a mystery to be explained by a
single parameter; it encodes the physical mass hierarchy
specific to each sector.

\paragraph{The shift hierarchy.}
The physical shifts form a hierarchy:
\begin{center}
\begin{tabular}{lcl}
$s \approx 1/\sqrt{6}$ & $\Rightarrow$ & $K = 2/3$ (leptons, inv.\ down quarks, $(t,b,c)$) \\
$s \approx 0.38$--$0.40$ & $\Rightarrow$ & $K \approx 0.46$--$0.66$ ($(b,c,s)$, $(c,s,u)$) \\
$s \approx 0.37$        & $\Rightarrow$ & $K \approx 0.73$ (down quarks, direct) \\
$s \approx 0.33$        & $\Rightarrow$ & $K \approx 0.85$ (up quarks) \\
$s \approx 0.49$        & $\Rightarrow$ & $K \approx 0.57$ ($(s,u,d)$) \\
\end{tabular}
\end{center}
The pattern $s_{\rm lep} \approx s_{(t,b,c)} \approx s_{(b,c,s)}
\approx z_0$ may reflect
the K\"ahler metric hierarchy:
composites involving only heavy quarks ($t$, $b$, $c$) have
K\"ahler metrics closer to canonical, while composites involving
light quarks ($u$, $d$, $s$) have larger corrections.

\paragraph{Predictive content.}
The free-shift ansatz with universal~$s$ has
\textbf{two} free parameters per sector:
$M$ (the overall scale) and $\alpha$ (the angle).
These determine three masses, yielding \textbf{one prediction}
per sector (the Koide ratio $K$ is predicted by the universal~$s$).
If $s = z_0$ is exact, this predicts $K = 2/3$ for all sectors;
deviations from~$2/3$ measure the K\"ahler correction $\delta_s$.

\paragraph{Relation to the quartic mixing.}
In the waterfall analysis~\cite{waterfall}, the inter-branch
coupling~$\varepsilon$ controls the mixing between the
direct and inverse branches.  The chain windows involve
both branches, and their effective shifts and angles are determined
by the quartic.  The universality of~$s$ in the chain
windows is therefore evidence that the quartic mixing
\emph{preserves} the $\UU(3)_F$ normalisation while
modifying the Cartan angle.

\paragraph{The Weyl duality.}
The $15^\circ/45^\circ$ split has a natural interpretation.
The $15^\circ$ fixed point is the ``standard'' Koide
configuration: one vanishing mass, the other two in the
golden ratio $R = (2+\sqrt{3})^2$.  The $45^\circ$ fixed
point is its Weyl mirror: one vanishing mass with a sign-flipped
charge.  Physical tuples near these fixed points
inherit the branch structure.  As a mass approaches zero
from the $15^\circ$ side, the charge $q = s + z$
passes through zero and the tuple transitions to
the $45^\circ$ branch.  The $(b,c,s)$ tuple, with its
near-zero strange-quark charge, is precisely at this
transition.

\paragraph{Open questions.}
\begin{enumerate}
\item \textbf{Origin of $s = 1/\sqrt{6}$.}
  The group-theoretic derivation of $s = z_0$ from
  $\UU(3)_F$ normalisation~\cite{letter_q2} explains the value
  but not why the physical K\"ahler potential is close to canonical.
  Can the $\sim 10\%$ sector-dependent corrections to~$s$ be
  computed from the K\"ahler metric of the confined $\SU(3)_F$ theory?

\item \textbf{Angle spectrum.}
  What determines the sector-dependent angles?
  The lepton angle ($12.7^\circ$) is close to the $15^\circ$
  fixed point, while the $(u,c,t)$ angle ($4.3^\circ$) is far
  from it.  In the M-theory lift, $\alpha$ is a K\"ahler modulus of
  the M5-brane curve.  Can the quark angles be predicted?

\item \textbf{Neutrino sector.}
  The baryonino masses (Majorana neutrinos) should also
  decompose into $(M, s, \alpha)$.  The universal~$s$
  predicts that the neutrino shift is the same
  as the charged-lepton shift, but the angle may differ.

\item \textbf{Weyl branch selection rule.}
  Is there a dynamical rule that determines whether a given
  mass window sits on the $15^\circ$ or $45^\circ$ branch?
  The empirical pattern is that the lightest mass in the
  window determines the branch: if its charge is positive, $15^\circ$;
  if negative (near the zero-mass boundary), $45^\circ$.
\end{enumerate}


%======================================================================
\section*{Acknowledgments}

During the preparation of this work the author used Claude
(Anthropic) for algebraic verification, numerical checks,
and drafting assistance.
The author reviewed and edited all content and takes full
responsibility for the publication.

\begin{thebibliography}{99}

\bibitem{Koide:1983qe}
  Y.~Koide,
  Phys.\ Lett.\ B \textbf{120}, 161 (1983).

\bibitem{PDG:2024}
  S.~Navas \textit{et al.} (Particle Data Group),
  Phys.\ Rev.\ D \textbf{110}, 030001 (2024).

\bibitem{Rivero:2005}
  A.~Rivero,
  arXiv:hep-ph/0505220 (2005).

\bibitem{Zenczykowski:2012}
  P.~Zenczykowski,
  Phys.\ Rev.\ D \textbf{86}, 117303 (2012).

\bibitem{letter_q2}
  A.~Rivero, ``Quadratic charge--mass relation from gauged
  family symmetry,'' companion letter (2025).

\bibitem{waterfall}
  A.~Rivero, ``The quark mass waterfall from
  $\mathcal{N}{=}2$ interleaving,'' companion letter (2025).

\bibitem{note98}
  A.~Rivero, ``Charge-plane geometry: zero-mass fixed points,''
  working note~98 (2025).

\bibitem{AHS:2005}
  A.~Rivero, ``AHS running Koide,'' working note~51 (2025);
  see also C.A.~Heusch, arXiv:0807.2094.

\bibitem{SMDR}
  S.P.~Martin and D.G.~Robertson,
  Phys.\ Rev.\ D \textbf{100}, 073004 (2019),
  arXiv:1907.02500.

\end{thebibliography}

\end{document}
